% =============================================================================
% Оршил (дугааргүй; «diplom ajluud»-ын Б.Болоржингийн ажлын бүтэцтэй ижил)
% =============================================================================

\defformat
\addchap{Оршил}
\label{ch:intro}
\pagecolor{white}

\addsec{Судалгааны үндэслэл, орчин}
\label{sec:motivation}

Банк бус санхүүгийн байгууллагын (ББСБ) өдөр тутмын шийдвэр тоо мэдээнд тулгуурладаг: салбар бүрийн хугацаа хэтэрсэн
зээлийн үлдэгдэл, сарын шинэ зээлийн тоо, барагдуулалтын үр дүн гэх мэт. Эдгээр мэдээллийг өгөгдлийн сангаас авахдаа
ажилтнууд SQL бичих чадвартай шинжээч эсвэл урьдчилан бэлтгэсэн тайланд хамаардаг. Шинэ асуулт гарах бүрд хүсэлт
илгээж хүлээх нь цаг алддаг бөгөөд шинжээчийн ачааллыг нэмэгдүүлдэг.

Байгалийн хэлний асуултыг SQL болгох (Text2SQL) системүүд том хэлний загварын (LLM) хөгжлийн хэрээр англи хэлний
жишгүүд дээр өндөр түвшинд хүрсэн \cite{shi_2025_survey_llm_text2sql, li_2023_bird}. Нээлттэй эхийн 7B хэмжээтэй
Arctic-Text2SQL-R1 загвар BIRD-ийн test дээр 68.5\% гүйцэтгэлийн нарийвчлал (EX) үзүүлдэг
\cite{yao_2025_arctic_text2sql_r1}. Гэвч англи бус хэл дээр эдгээр загварын нарийвчлал мэдэгдэхүйц буурдаг
\cite{dou_2023_multispider, pham_2026_multispider2}. Санхүүгийн өгөгдөл байгууллагаас гадагш гарах ёсгүй тул арилжааны
API-д суурилсан шийдэл тохиромжгүй бөгөөд нэг GPU-д багтах нээлттэй эхийн загварыг монгол хэл ба байгууллагын домэйнд
тохируулах шаардлагатай.

\addsec{Асуудлын тодорхойлолт}
\label{sec:problem}

Монгол асуулт $q$, өгөгдлийн сангийн схем $S$ ба байгууллагын мэдлэг $K$ (баганын тайлбар, нэр томьёо, жишээ)
өгөгдсөн үед гүйцэтгэлийн үр дүн нь лавлах (gold) SQL-ийнхтэй таарах SQL $\hat{s}$-ийг үүсгэх, эсвэл асуулт хариулах
боломжгүй (хувийн мэдээлэл, өгөгдөлд байхгүй) бол татгалзах, хоёрдмол бол тодруулга асуух шаардлагатай. Энэ даалгаврыг
дараах бэрхшээлүүд хүндрүүлнэ:
\begin{enumerate}
  \item \textbf{Бага нөөцтэй хэл.} Монгол Text2SQL-ийн сургалтын ба үнэлгээний өгөгдөл байхгүй, загварууд монгол
        хэлийг сул ойлгодог \cite{zhang_2024_mm_eval}. Хэрэглэгчид кирилл ба латин галигаар («zeeliin hvselt») холимог бичдэг.
  \item \textbf{Домэйны мэдлэг.} «PAR30», «идэвхтэй багц» зэрэг нэр томьёоны тодорхойлолт байгууллагын дотоод дүрмээс
        хамаардаг \cite{li_2023_bird}.
  \item \textbf{Том схем.} 28 view, 425 багана бүхий схемийг бүтнээр нь өгөхөд prompt 13--20 мянган токен болж, жижиг
        загварын анхаарлыг сарниулдаг.
  \item \textbf{Найдвартай байдал.} Санхүүгийн хэрэглээнд буруу тоо хариулахаас татгалзах нь дээр
        \cite{lee_2024_trustsql}.
\end{enumerate}

\addsec{Таамаглал}
\label{sec:hypothesis}

\textbf{H1:} Мэдлэг хайлт (схемийн тайлбар, утга, нэр томьёо, жишээ) нь нийт зөв хариултыг 10-аас дээш пунктээр өсгөнө.

\textbf{H2:} Англи төвтэй SQL загварт монгол асуултыг англи руу орчуулж өгөх нь шууд өгөхөөс статистикийн ач
холбогдолтой ($p < 0.05$) илүү үр дүн өгнө.

\textbf{H3:} Өөрийн өгөгдлийн сангийн RAG prompt-оор нэмж сургасан (fine-tune) загвар RAG-тай суурь загварыг давна.

\textbf{H4:} Олон хувилбараас гүйцэтгэлийн үр дүнгээр санал хураах нь хариулах боломжтой асуултын нарийвчлалыг ач
холбогдолтой өсгөнө.

\addsec{Зорилго}
\label{sec:goal}

ББСБ-ийн өгөгдлийн сан дээр монгол хэлээр асуусан асуултыг автоматаар SQL болгож хариулдаг систем бүтээх.

\addsec{Зорилтууд}

Дээрх зорилгод хүрэхийн тулд сэдэв батлах даалгаврын дагуу дараах зорилтуудыг дэвшүүлэв:
\begin{enumerate}
  \item Өгөгдлийн сангийн схем, утга, бизнесийн нэр томьёо, жишээ асуултаас мэдлэгийн сан (RAG) бүрдүүлэх;
  \item Асуулт бүрд хамааралтай мэдээллийг хайж олоод загварт өгөх prompt болон pipeline боловсруулах;
  \item Загварыг монгол хэл болон өөрийн өгөгдлийн санд тохируулан сургах (fine-tune);
  \item Хариултын зөв байдлыг нэмэгдүүлэх: олон хувилбараас сонгох, алдааг засах, хариулах боломжгүй асуултаас татгалзах;
  \item Системийн үр дүнг тестийн асуултын багцаар хэмжиж үнэлэх, тайлагнах.
\end{enumerate}

\addsec{Хамрах хүрээ ба хязгаарлалт}
\label{sec:scope}

\textbf{Хамрах хүрээ:}
\begin{itemize}
  \item ББСБ-ийн PostgreSQL өгөгдлийн сан (28 view бүхий text2sql давхарга) дээрх монгол асуултад SQL үүсгэх систем;
  \item Кирилл болон латин галигаар бичсэн асуулт, хариулах боломжгүй ба хоёрдмол асуултыг таних;
  \item Нэг GPU-д ажиллах нээлттэй эхийн загвар (SQL бичигчээр зөвхөн Arctic-Text2SQL-R1-7B).
\end{itemize}

\textbf{Хязгаарлалт:}
\begin{itemize}
  \item Бодит харилцагчийн мэдээлэл ашиглах эрхгүй тул детерминист синтетик өгөгдөл (seed 42, 2026-06-30-ны байдлаар) ашиглав;
  \item Өөрийн GPU сервергүй тул үнэлгээг Kaggle T4×2 (долоо хоногт 30 цаг), сургалтыг Modal ба Vast.ai-ийн A100 дээр
        хязгаарлагдмал төсвөөр хийсэн; иймд төлөвлөгөөний зарим туршилт (GRPO, agentic горим) хийгдээгүй;
  \item Гол хэмжилт кирилл асуулт дээр хийгдсэн; бодит хэрэглээний (production) баталгаажуулалт хамрагдаагүй.
\end{itemize}

\addsec{Судалгааны хувь нэмэр}
\label{sec:contrib-intro}

\begin{enumerate}
  \item Хариулах боломжгүй ба хоёрдмол асуулттай, кирилл, латин, англи хувилбартай, загварын гэр бүлээр хуваасан 500
        асуулт бүхий монгол хэлний Text2SQL үнэлгээний сан;
  \item Мэдлэг хайлт, орчуулга, нэмэлт сургалт, санал хураалт тус бүрийн нөлөөг нэг схем, нэг үнэлгээн дээр McNemar
        тестээр харьцуулсан эмпирик нотолгоо;
  \item Алдааны эх үүсвэрийн хэмжилт ба түүнд тулгуурласан санхүүгийн хэллэгийн орчуулгын өгөгдөл (2 828 хос);
  \item Үнэгүй үүлэн квот дээр тасалдлаас сэргэдэг туршилтын систем, алхам бүрийн prompt ба мэдлэгийг харуулдаг
        «Text2SQL Studio» хэрэгсэл.
\end{enumerate}

\addsec{Ажлын бүтэц}
\label{sec:structure}

Бүлэг~\ref{ch:related}-т холбоотой судалгаа ба ижил төстэй бүтээгдэхүүнийг тоймлож судалгааны цоорхойг тодорхойлно.
Бүлэг~\ref{ch:method}-т санал болгож буй архитектур, SQL бичигч ба орчуулгын загварын онолын үндсийг, Бүлэг~\ref{ch:data}-т
өгөгдлийн сан, мэдлэгийн сан, сургалтын ба үнэлгээний өгөгдлийг, Бүлэг~\ref{ch:model}-т загваруудын сургалтын урсгал ба
туршилтын тохиргоог тайлбарлана. Бүлэг~\ref{ch:system}-д туршилтын системийн хэрэгжүүлэлтийг, Бүлэг~\ref{ch:results}-д үр
дүн, хэлэлцүүлэг ба ёс зүйн асуудлыг авч үзээд, эцэст нь Дүгнэлт хэсэгт нэгтгэн дүгнэнэ.

\myformat
